\section{ИССЛЕДОВАНИЕ И ОЦЕНКА КАЧЕСТВА РАЗРАБОТАННОГО РЕШЕНИЯ}

После реализации компонентов платформы необходимо проверить, что полученное решение соответствует поставленной задаче и корректно выполняет ключевой цикл управления: построение desired state, получение actual state, обнаружение drift и запуск reconcile. Оценка качества в данной работе выполняется по двум направлениям. Первое направление связано с тестированием отдельных сервисов и их доменной логики. Второе направление связано с экспериментальной оценкой поведения платформы на сценариях, отражающих основные эксплуатационные свойства: масштабируемость detection cycle, время сходимости после drift, пропускную способность reconcile, устойчивость к partial data и эффективность cooldown.

\subsection{Цели, критерии и ограничения испытаний}

Главная цель испытаний --- подтвердить, что разработанная платформа реализует заявленную архитектурную модель. Для этого необходимо проверить, что каждый сервис корректно выполняет собственную функцию, а цепочка сервисов в совокупности обеспечивает обнаружение и устранение поддержанного drift.

Критерии оценки сформулированы следующим образом:
\begin{enumerate}
    \item ParserSvc должен корректно загружать Structurizr JSON, валидировать данные и строить desired state;
    \item InventorySvc должен корректно формировать actual state, поддерживать snapshot и явно отражать partial result;
    \item DriftDetectorSvc должен корректно сопоставлять desired state и actual state, обнаруживать drift только при достаточных данных и подавлять повторные reconcile-команды в пределах cooldown;
    \item ReconcilerSvc должен принимать reconcile-команды, валидировать payload, выбирать operator, строить execution plan и ставить задачу в очередь;
    \item API-контракты сервисов должны быть согласованы с фактическими структурами данных;
    \item экспериментальные сценарии должны демонстрировать работоспособность control loop и позволять оценить узкие места MVP.
\end{enumerate}

Испытания имеют несколько ограничений. Во-первых, текущая версия платформы поддерживает только два компонента: \texttt{node\_exporter} и \texttt{cadvisor}. Следовательно, результаты нельзя напрямую переносить на произвольные workload без реализации дополнительных detector-ов и operator-ов. Во-вторых, состояние cooldown, desired state и actual state хранится в памяти, поэтому верификация не охватывает сценарии восстановления после перезапуска с сохранением истории. В-третьих, измерения производительности не являются SLA: они отражают поведение MVP в выбранных условиях и служат инженерной точкой отсчета для последующих оптимизаций.

Отдельно важно зафиксировать ограничение partial data. В платформе отсутствие фактических данных не должно автоматически считаться drift. Поэтому в испытаниях проверяется не только способность обнаруживать расхождения, но и способность не выполнять лишние reconcile-действия при неполной информации.

\subsection{Модульные и компонентные тесты}

Тестовый контур охватывает все основные сервисы платформы. Для Go-сервисов используются unit-тесты, расположенные рядом с соответствующими пакетами. Для ReconcilerSvc используются Python-тесты на базе pytest. Такой подход соответствует структуре проекта: каждый сервис проверяется в собственной технологической среде и на уровне собственных доменных компонентов.

В ParserSvc тестируются загрузка конфигурации, обнаружение manifest-файлов, разбор Structurizr JSON, mapping deployment-модели в desired state, валидация workload и host-полей, in-memory store и HTTP handlers. Эти тесты подтверждают, что сервис корректно обрабатывает валидные манифесты, пропускает некорректные элементы и поддерживает фильтрацию desired-хостов.

В InventorySvc проверяются загрузка конфигурации, валидация selfProvisioning, mapping ответа cAdvisor, фильтрация системных workload, потокобезопасное поведение in-memory store, фильтрация хостов по labels и HTTP handlers. Особое значение имеют тесты cAdvisor mapper: они подтверждают, что сервис выделяет из внешнего ответа именно workload, которые затем могут использоваться DriftDetectorSvc для сравнения с desired state.

В DriftDetectorSvc тестируются конфигурация, HTTP-клиенты ParserSvc, InventorySvc и ReconcilerSvc, detector registry, detector-ы для \texttt{node\_exporter} и \texttt{cadvisor}, cooldown store, scheduler, HTTP handlers и основной detection service. Набор тестов включает проверку обнаружения отсутствующего компонента, отсутствие drift при наличии workload, bootstrap-поведение для \texttt{cadvisor}, подавление повторных reconcile-команд, partial-обработку при отсутствующем actual host и запрет параллельного запуска detection cycle.

В ReconcilerSvc тестируются схемы Pydantic, API endpoints, operator registry, построение execution plan, очередь задач, reconcile service и Ansible adapter с подменой runner-а. Эти тесты подтверждают, что сервис корректно валидирует входные данные, возвращает \texttt{202 Accepted} при успешной постановке в очередь, возвращает \texttt{503} при насыщении очереди и строит корректный in-memory inventory для Ansible Runner.

\begin{table}
    \centering
    \caption{Основные группы тестов компонентов платформы}
    \label{tab:ch4-test-groups}
    \footnotesize
    \begin{tabularx}{\textwidth}{|p{0.22\textwidth}|X|}\hline
    Сервис & Проверяемые аспекты \\ \hline
    ParserSvc & Конфигурация, поиск JSON-манифестов, разбор Structurizr JSON, mapping desired state, валидация, store, API handlers \\ \hline
    InventorySvc & Конфигурация, selfProvisioning, cAdvisor client/mapper, фильтрация workload, in-memory snapshot, labels filtering, API handlers \\ \hline
    DriftDetectorSvc & HTTP-клиенты, detector registry, detector-ы, cooldown, detection cycle, scheduler, partial-обработка, API handlers \\ \hline
    ReconcilerSvc & Pydantic-схемы, API, operator registry, execution plan, очередь, reconcile service, Ansible adapter \\ \hline
    \end{tabularx}
    \normalsize
\end{table}

Таким образом, модульные и компонентные тесты покрывают наиболее рискованные участки реализации: преобразование данных между форматами, принятие решения о drift, подавление повторных действий и асинхронную постановку reconcile-задач. Это создает основу для дальнейшей экспериментальной оценки поведения платформы как единой системы.

\subsection{Методика экспериментальной оценки}

Экспериментальная оценка выполнялась набором сценариев, которые проверяют не отдельные функции сервисов, а поведение платформы как связанного control loop. Для этого использовались HTTP API сервисов, структурированные логи и сохраненные артефакты выполнения.

Методика включает пять сценариев:
\begin{enumerate}
    \item оценка масштабируемости detection cycle при различном числе управляемых хостов;
    \item оценка времени сходимости после внесения drift;
    \item оценка пропускной способности ReconcilerSvc и поведения при насыщении очереди;
    \item проверка устойчивости при partial data;
    \item оценка эффективности cooldown для подавления повторных reconcile-команд.
\end{enumerate}

Для каждого сценария собирались сырые результаты в JSON и сводный отчет. В качестве метрик использовались длительность detection cycle, длительности отдельных стадий цикла, количество найденных drift, количество отправленных и подавленных reconcile-команд, число принятых и отклоненных reconcile-запросов, throughput завершения reconcile-задач и признаки partial-обработки.

Важная особенность методики состоит в разделении стадий detection cycle. Для анализа недостаточно знать только полную длительность цикла. Необходимо понимать, какая часть времени приходится на получение inventory, какая --- на получение desired state, какая --- на сравнение, а какая --- на dispatch reconcile-команд. Поэтому DriftDetectorSvc сохраняет stage timings: \texttt{InventoryFetchMs}, \texttt{ParserFetchMs}, \texttt{DriftComparisonMs}, \texttt{ReconcileDispatchMs}.

\subsection{Масштабируемость detection cycle}

Первый сценарий направлен на оценку того, как длительность detection cycle изменяется при росте числа управляемых хостов. В сценарии рассматривались случаи с 1, 10, 100 и 500 хостами. Для каждого случая выполнялось несколько циклов в устойчивом состоянии, когда reconcile dispatch отсутствовал. Это позволяет оценивать именно стоимость получения и сравнения состояний, а не стоимость восстановления workload.

\begin{table}
    \centering
    \caption{Результаты оценки масштабируемости detection cycle}
    \label{tab:ch4-detection-scalability}
    \footnotesize
    \begin{tabular}{|r|r|r|r|r|r|r|}\hline
    Хосты & Cycle p50, мс & Inventory p50, мс & Parser p50, мс & Compare p50, мс & Dispatch p50, мс \\ \hline
    1   & 0.0  & 0.0  & 0.0 & 0.0 & 0.0 \\ \hline
    10  & 1.0  & 0.0  & 0.0 & 0.0 & 0.0 \\ \hline
    100 & 6.0  & 3.0  & 2.0 & 0.0 & 0.0 \\ \hline
    500 & 20.0 & 12.0 & 6.0 & 0.0 & 0.0 \\ \hline
    \end{tabular}
    \normalsize
\end{table}

Данные таблицы~\ref{tab:ch4-detection-scalability} показывают, что рост длительности цикла в основном связан с увеличением времени получения actual state и desired state. Для 100 хостов медианная длительность цикла составила 6 мс, для 500 хостов --- 20 мс. При этом медианное время стадии сравнения округлялось к 0 мс, что показывает низкую стоимость самой логики сопоставления в сравнении с fetch-стадиями.

Отсутствие dispatch-времени объясняется тем, что измеренные циклы выполнялись в состоянии без новых reconcile-команд. Это важно для корректной интерпретации: сценарий оценивает масштабирование detection logic, а не производительность ReconcilerSvc.

Полученный результат подтверждает, что в текущей реализации основная стоимость detection cycle при росте числа хостов связана с получением данных от сервисов-источников. Следовательно, дальнейшие оптимизации масштабируемости должны быть направлены прежде всего на ускорение InventorySvc/ParserSvc API, кэширование, уменьшение объема передаваемых данных или введение более точных инкрементальных механизмов.

\subsection{Время сходимости после drift}

Второй сценарий оценивает время сходимости после искусственного появления drift. Под сходимостью понимается интервал от момента внесения расхождения до момента восстановления требуемого workload и последующего исчезновения drift в наблюдаемом состоянии.

В измеренном сценарии удалялся workload \texttt{node\_exporter}. Далее фиксировались ключевые интервалы: время до первого detection cycle, время от обнаружения до отправки reconcile-команды, время выполнения восстановления и полное время сходимости.

\begin{table}
    \centering
    \caption{Время сходимости после drift}
    \label{tab:ch4-convergence-time}
    \footnotesize
    \begin{tabularx}{0.8\textwidth}{|X|r|}\hline
    Метрика & Значение, мс \\ \hline
    Inject -> first detection & 4075 \\ \hline
    Detection -> reconcile send & 4 \\ \hline
    Reconcile send -> workload restored & 10219 \\ \hline
    Total convergence & 14300 \\ \hline
    \end{tabularx}
    \normalsize
\end{table}

Как видно из таблицы~\ref{tab:ch4-convergence-time}, полное время сходимости составило 14300 мс. При этом собственно отправка reconcile-команды после обнаружения drift заняла 4 мс. Основная часть времени пришлась на выполнение восстановления и ожидание обновления actual state.

Такой результат согласуется с архитектурой платформы. DriftDetectorSvc быстро принимает решение и передает команду ReconcilerSvc, но фактическое восстановление зависит от выполнения Ansible playbook и последующего появления обновленных данных в inventory snapshot. Следовательно, время сходимости определяется не только скоростью detection cycle, но и длительностью reconcile execution, частотой обновления inventory и задержками источников actual state.

Практический вывод состоит в том, что для сокращения convergence latency необходимо оптимизировать несколько участков одновременно. Уменьшение интервала detection само по себе не даст полного эффекта, если reconcile playbook выполняется долго или actual state обновляется редко. В дальнейшем полезно добавить явные события о завершении reconcile и более точное отслеживание момента восстановления workload.

\subsection{Throughput reconcile и насыщение очереди}

Третий сценарий направлен на проверку ReconcilerSvc. В нем оценивалось, как меняется throughput завершения reconcile-задач при разных значениях параллелизма, а также как сервис ведет себя при burst-нагрузке.

\begin{table}
    \centering
    \caption{Результаты оценки reconcile throughput}
    \label{tab:ch4-reconcile-throughput}
    \footnotesize
    \begin{tabular}{|p{0.28\textwidth}|r|r|r|r|r|r|}\hline
    Сценарий & Parallelism & Requests & Accepted & Rejected & Success & Ops/s \\ \hline
    steady\_parallel\_1 & 1 & 90  & 90  & 0  & 90 & 0.2392 \\ \hline
    steady\_parallel\_2 & 2 & 90  & 90  & 0  & 81 & 0.4636 \\ \hline
    steady\_parallel\_4 & 4 & 90  & 90  & 0  & 83 & 0.7822 \\ \hline
    saturation\_probe\_parallel\_2 & 2 & 180 & 102 & 78 & 99 & 0.4579 \\ \hline
    \end{tabular}
    \normalsize
\end{table}

Результаты таблицы~\ref{tab:ch4-reconcile-throughput} показывают ожидаемую зависимость: при увеличении \texttt{max\_parallel\_reconciliations} throughput завершения возрастает.

Отдельный интерес представляет сценарий saturation. При 180 запросах и параллелизме 2 сервис принял 102 запроса и отклонил 78. Это подтверждает, что ограничение очереди работает: ReconcilerSvc не принимает неограниченное число задач, а возвращает \texttt{503}, когда очередь не может принять новую задачу в заданный timeout. С точки зрения надежности это предпочтительнее бесконтрольного накопления задач и роста задержек.

Следует учитывать, что число successful completions в некоторых сценариях меньше числа accepted. Это отражает асинхронную природу выполнения и ограниченное окно наблюдения за завершением задач. Для более точной production-оценки в дальнейшем требуется persisted job state и отдельный API статуса очереди, чтобы различать принятые, выполняемые, завершенные и ошибочные задачи.

Таким образом, текущая реализация ReconcilerSvc демонстрирует управляемое поведение под нагрузкой. Worker pool позволяет масштабировать параллелизм, а queue capacity задает верхнюю границу накопления задач. Главным узким местом остается выполнение Ansible playbook, что ожидаемо для agentless-модели восстановления через SSH.

\subsection{Устойчивость к partial data и эффективность cooldown}

Четвертый и пятый сценарии проверяют безопасность автоматических действий. Для инфраструктурной платформы важно не только уметь устранять drift, но и не выполнять лишние действия при неполных данных или повторяющихся сигналах.

В сценарии partial data рассматривалась ситуация, когда часть хостов недоступна для получения actual state. В измерении использовалось 20 хостов, из которых 6 имели недоступные FQDN для InventorySvc. Результаты приведены в таблице~\ref{tab:ch4-partial-data}.

\begin{table}
    \centering
    \caption{Результаты проверки устойчивости к partial data}
    \label{tab:ch4-partial-data}
    \footnotesize
    \begin{tabularx}{0.85\textwidth}{|X|r|}\hline
    Метрика & Значение \\ \hline
    Cycle partial flag & true \\ \hline
    Inventory \texttt{metadata.isPartial} & true \\ \hline
    Inventory \texttt{failedHosts} & 6 \\ \hline
    Warnings: missing actual data & 6 \\ \hline
    Node exporter reconciles for missing hosts & 0 \\ \hline
    Cadvisor reconciles for missing hosts & 0 \\ \hline
    \end{tabularx}
    \normalsize
\end{table}

Полученные значения подтверждают корректную partial-семантику. DriftDetectorSvc продолжил обработку доступных данных, но не стал автоматически отправлять reconcile-команды для хостов с недостоверным actual state. Для отсутствующих данных были зафиксированы предупреждения, а цикл был помечен как partial. Это соответствует требованию безопасности: неполнота данных не должна превращаться в основание для произвольного восстановления.

Сценарий cooldown проверял подавление повторных reconcile-команд. Он моделировал повторяющийся drift в нескольких циклах. Результаты приведены в таблице~\ref{tab:ch4-cooldown}.

\begin{table}
    \centering
    \caption{Результаты проверки эффективности cooldown}
    \label{tab:ch4-cooldown}
    \footnotesize
    \begin{tabularx}{0.85\textwidth}{|X|r|}\hline
    Метрика & Значение \\ \hline
    Reconcile sent without cooldown & 20 \\ \hline
    Reconcile sent with cooldown & 0 \\ \hline
    Prevented reconcile commands & 20 \\ \hline
    Suppression ratio & 1.0 \\ \hline
    \end{tabularx}
    \normalsize
\end{table}

Результаты показывают, что cooldown полностью подавил повторные отправки в рамках окна эксперимента. Это важно для устойчивости платформы. Если drift не может быть устранен сразу, повторная отправка одинаковых команд в каждом цикле привела бы к перегрузке ReconcilerSvc и управляемых хостов. Cooldown снижает этот риск и делает поведение системы более предсказуемым.

Вместе с тем текущая реализация cooldown хранит состояние в памяти DriftDetectorSvc. Следовательно, при перезапуске сервиса история подавления теряется. Для более зрелой эксплуатации необходимо вынести cooldown/history в персистентное хранилище или отдельный механизм deduplication.

\subsection{Интерпретация результатов и ограничения измерений}

Проведенные испытания показывают, что разработанная платформа выполняет ключевые функции MVP. ParserSvc корректно строит desired state из архитектурного JSON-представления, InventorySvc формирует actual state и явно отражает partial result, DriftDetectorSvc обнаруживает поддержанные виды drift и подавляет повторные reconcile-команды, а ReconcilerSvc принимает команды и выполняет их через асинхронную очередь и operator-слой.

Экспериментальная оценка подтверждает несколько важных свойств. Во-первых, detection cycle масштабируется предсказуемо для проверенного диапазона хостов, а основная стоимость роста связана с получением данных от сервисов-источников. Во-вторых, время сходимости после drift определяется не только скоростью обнаружения, но и длительностью reconcile и обновления actual state. В-третьих, ReconcilerSvc демонстрирует управляемое поведение под нагрузкой: рост параллелизма повышает throughput, а насыщение очереди приводит к отказу в приеме новых задач вместо неконтролируемого накопления. В-четвертых, partial data и cooldown работают как защитные механизмы от ошибочных или избыточных действий.

Ограничения измерений также существенны. Сценарии с 100 и 500 хостами отражают масштабирование control-plane-логики, а не полноценную эксплуатацию сотен независимых физических узлов. Абсолютные значения задержек зависят от вычислительных ресурсов, конфигурации сервисов, особенностей Ansible-выполнения и способа моделирования данных. Кроме того, stage timing для очень быстрых стадий может округляться до 0 мс, поэтому такие значения следует интерпретировать как признак малой длительности, а не как буквальное отсутствие затрат.

На основе результатов можно выделить направления улучшения качества. Для detection cycle полезно добавить более подробные метрики по стадиям и объемам данных. Для ReconcilerSvc необходим API статуса очереди и persisted job history. Для cooldown и истории drift требуется персистентное состояние. Для convergence-оценки полезен явный event о восстановлении workload, чтобы уменьшить зависимость измерений от периодического получения actual state.

Таким образом, глава подтверждает, что разработанное решение является работоспособным MVP автоматизированной платформы управления инфраструктурой с поддержкой Architecture-as-Code. Оно реализует основной control loop и демонстрирует ожидаемое поведение в ключевых сценариях, при этом имеет понятные ограничения и направления дальнейшего развития.
